% concatenated from phys.tex
\documentclass[reprint,amsmath,amssymb,aps,]{revtex4-2}

\usepackage[T1]{fontenc}
\usepackage{lmodern}
\usepackage{textcomp}
\usepackage{latexsym}                                                               
\usepackage{amsfonts}
\usepackage[dvipdfmx]{graphicx}
\usepackage{amsthm,amsmath,amssymb,bm}            
\usepackage{fancybox,url,ascmac,comment,wrapfig}
\usepackage{mathtools}
\usepackage{graphicx}
\usepackage{ascmac}
\usepackage{physics}
\usepackage{tikz}
\usepackage[inline]{enumitem} 
\usepackage{algorithmic}
\usepackage{algorithm}

\theoremstyle{definition}
\newtheorem{exer}{Exercise}
\newtheorem{defi}{Definition}
\newtheorem{exam}{Example}
\newtheorem{lem}{Lemma}
\newtheorem{theo}{Theorem}
\newtheorem{note}{Note}
\newtheorem{prop}{Proposition}
\newtheorem{coro}{Corollary}
\newtheorem{rem}{Remark}
\newtheorem{fact}{Fact}
\newtheorem{algo}{Algorithm}
\newtheorem{proc}{Procedure}

%%色文字
\newcommand{\R}[1]{\textcolor{red}{#1}}
\colorlet{RED}{red}
\newcommand{\Corr}[1]{\textcolor{red}{#1}}
\newcommand{\B}[1]{\textcolor{blue}{#1}}

%%白抜き文字
\newcommand{\SET}[1]{\mathbb{#1}}
\newcommand{\NNR}{\SET{R}_{\ge 0}}
\newcommand{\Z}{\SET{Z}}
\newcommand{\N}{\SET{N}}
\newcommand{\C}{\SET{C}}
\newcommand{\1}{\mbox{1}\hspace{-0.25em}\mbox{l}}%%指示関数

%%整数集合
\newcommand{\setint}[1]{[\hspace{-0.5mm}[#1]\hspace{-0.5mm}]}
\newcommand{\subint}[2]{\tbinom{\setint{#1}}{#2}}

%行列
\newcommand{\mat}[1]{\mathbf{#1}}

%密度行列の集合
\newcommand{\pure}[2]{\C^{#1 \otimes #2}}
\newcommand{\state}[2]{S(\pure{#1}{#2})}
\newcommand{\QU}[1]{\mathcal{#1}}

%% \affiliate[Time]{The author is with The Research Institute for Time Studies, Yamaguchi University}

\begin{document}
\preprint{APS/123-QED}

\title{Insertion Correcting Capability for Quantum Deletion-Correcting Codes}

\author{Ken Nakamura}
\email{f006wbw@yamaguchi-u.ac.jp}
\author{Takayuki Nozaki}%
\email{tnozaki@yamaguchi-u.ac.jp}
\affiliation{
  %% Authors' institution and/or address\\
  %% This line break forced with \textbackslash\textbackslash
  Dept.\ of Informatics, Division of Fundamental Sciences,
  Graduate School of Sciences and Technology for Innovation, 
  Yamaguchi University,
  1677-1, Yoshida, Yamaguchi-shi, Yamaguchi, 753-8512, JAPAN
}
\altaffiliation[Also at]{
  the Research Institute for Time Studies, Yamaguchi University
}

\date{\today}% It is always \today, today,
             %  but any date may be explicitly specified

\begin{abstract}
  This paper proves that
  any quantum $t$-deletion-correcting codes
  also correct a total of $t$ insertion and deletion errors
  under a certain condition.
  Here,
  this condition
  is that
  a set of quantum states is defined as a quantum error-correcting code
  if
  the error spheres of its states are disjoint,
  as classical coding theory.
  In addition,
  this paper proposes the quantum indel distance
  and
  describes insertion and deletion errors correcting capability of quantum codes
  by this distance.
\end{abstract}

\maketitle

\section{Introduction}
In classical coding theory,
a classical insertion error adds a new symbol to a sequence
and
a classical deletion error loses a symbol of a sequence.
A set of sequences
after
classical deletion/insertion errors to a sequence $\bm{c}$
is called
a {\it deletion/insertion sphere centered at $\bm{c}$.}
Classical deletion/insertion-correcting codes
are defined by
the sets of sequences
such that
deletion/insertion spheres centered at their elements
are disjoint.
It is known that
a code $C$ corrects $t$ classical deletion errors
if and only if
$C$ corrects a total of $t$ classical deletion and insertion errors
\cite{levenshtein1966binary}.

In quantum coding theory,
deletion and insertion errors are defined as
classical coding theory;
namely,
an insertion error adds a new qudit to a multi-qudit system
and
a deletion error loses a qudit of a multi-qudit system.
In previous research,
the quantum deletion-correcting codes
\cite{hagiwara2020four,matsumoto2022constructions,hagiwara2023quantum}
correct insertion errors
\cite{hagiwara2021four,nakamura2025multiple,sasaki2024insertion}.
In other words,
these quantum deletion-correcting codes
are also
quantum insertion-correcting codes.
However,
it is an open problem \cite{matsumoto2021survey,hagiwara2024introduction}
that
for a quantum code $\QU{C}$,
the following statements are equivalent:
\begin{enumerate*}[label={(\roman*)}]
\item
  $\QU{C}$
  corrects
  $t$ deletion errors.
\item
  $\QU{C}$
  corrects
  a total of $t$ deletion and insertion errors.
\end{enumerate*}

To solve this problem,
Shibayama et al. 
focused on the Knill-Laflamme (KL) condition.
They 
showed that
KL condition for $t$ deletion errors and
that for $t$ {\it separable insertion errors}
are 
equivalent \cite{shibayama2022equivalence}.
Here,
a separable insertion error is
an insertion error where
there is no entanglement
between the original multi-qudit system and the inserted qudit.
It is known that
insertion errors
are limited to
separable insertion errors
if
the original multi-qudit system is pure \cite{sibayama_single_ins_def}.
Furthermore,
a quantum code $\QU{C}$ satisfies
the KL condition for deletion/insertion errors
if and only if
there is
a recovery operator of $\QU{C}$ \cite{shibayama2025necessary}.
In other words,
Shibayama et al. \cite{shibayama2022equivalence,shibayama2025necessary} proved that
a quantum code $\QU{C}$ corrects
$t$ deletion errors
if and only if
$\QU{C}$ corrects a total of $t$ deletion and separable insertion errors.
In particular, 
a quantum code $\QU{C}$ corrects
single-deletion errors
if and only if
$\QU{C}$ corrects single-insertion errors.

There
are
two issues
in \cite{shibayama2022equivalence,shibayama2025necessary}.
A first issue
is
to limit insertion errors to separable insertion errors.
In the case of
composite errors consisting of deletions and insertions,
any insertion error is not necessarily limited to
separable insertion error
since deletion/insertion errors
may change to mixed states from pure states.
Hence,
the discussion in \cite{shibayama2025necessary}
was limited to single-deletion/insertion errors
and
this result is not applied to composite errors.
Furthermore,
separable insertion errors
are
strong assumptions
since
deletion and insertion errors are non-commutative \cite{hagiwara2024introduction},
deletion and separable insertion errors are not.
A second issue is that
the KL condition for deletion/insertion errors in \cite{shibayama2022equivalence,shibayama2025necessary}
is not 
suitable
to discuss
the deletion/insertion correcting capability.
More precisely,
this KL condition
also requires 
the orthogonality of the basis after deletion/insertion errors
for distinct error positions.
Indeed,
there exist
a quantum code \cite{hagiwara2024pure}
such that
the basis after single-deletion error
is
orthogonal
but
distinct codewords
change
to
the same state
by this error;
namely,
this code
satisfies 
the orthogonality of the basis after single-deletion
but
cannot correct this error.
From the above,
the condition for the deletion/insertion correcting capability
also requires that 
distinct codewords
do not change
to
the same state
by deletion/insertion errors.

To solve the open problem and these issues,
we focus on the deletion/insertion sphere.
Here,
the deletion/insertion sphere
is defined as classical coding theory;
namely,
the deletion/insertion sphere centered at a quantum state $\rho$
is
a set of quantum states
after
deletion/insertion errors to $\rho$.

The main theorem (Theorem \ref{theo:1del_iff_1ins_2}) in this paper
is that
any quantum $t$-deletion-correcting codes
also correct a total of $t$ insertion and deletion errors,
where
a quantum deletion/insertion-correcting code
is defined
by
a set of quantum states
such that 
the deletion/insertion spheres of its states are disjoint,
similar to classical coding theory.
The main theorem is derived
from a proposition (Theorem \ref{theo:insdel_sphere}) that
any composite errors consisting of $t$-insertion and $s$-deletion errors
are in
an error where $t$-insertion errors occur after $s$-deletion errors.

In addition,
this paper explicitly provides
quantum states after insertion error to mixed states.
The previous research \cite{sibayama_single_ins_def} 
explicitly provided
quantum states after insertion error to pure states.
This result is not applied
in cases of composite errors consisting of insertions and deletions
since 
deletion/insertion errors
may change
to mixed states
from pure states.
By generalizing the result in \cite{sibayama_single_ins_def},
this paper explicitly provides
quantum states after insertion error to mixed states.
From this,
we can calculate
quantum states after insertion error
and
the set of its quantum state.
By using this result,
this paper gives
\begin{enumerate*}[label={(\roman*)}]
\item an example of the main theorem and Theorem \ref{theo:insdel_sphere},
\item a set of quantum states
  such that
  the set corrects single-insertion
  but
  not a single-deletion.
\end{enumerate*}
This
suggests that
the converse of the main theorem
does not
hold,
unlike in classical coding theory.
In other words,
any quantum insertion-correcting code
does not necessarily
correct
deletion errors.

From the main theorem and Theorem \ref{theo:insdel_sphere},
this paper
defines
the quantum indel distance
and
characterizes the insertion and deletion errors correcting capability of quantum codes
by this distance. 
The Hamming distance \cite{hamming1950error},
the indel distance \cite{levenshtein1966binary}, and
the distance of a stabilizer code \cite{nielsen2016quantum}
characterize
the error correcting capability
for
bit-flip errors,
classical insertion/deletion errors,
and unitary errors,
respectively,
and
derive
the bounds of classical/quantum error-correcting codes
for each error.
To characterize insertion and deletion errors correcting capability
and
derive some bound of quantum deletion/insertion-correcting codes,
this paper
defines
the quantum indel distance
based on the indel distance \cite{levenshtein1966binary}.
By this distance,
this paper
proves that
a set $\QU{X}$ 
corrects
a total of $t$ insert
ion and deletion errors
if
the quantum indel distance of distinct states in $\QU{X}$ 
is
greater than $2t$.

To summarize the above,
the contributions of this paper
are
as follows:
\begin{itemize}
\item
  This paper
  proves that
  any quantum $t$-deletion-correcting codes
  also correct a total of $t$ insertion and deletion errors,
  where
  a set of quantum states is defined as a quantum error-correcting code
  if
  the error spheres of its states are disjoint.
\item  
  This paper explicitly provides
  quantum states after insertion error to mixed states.
\item
  This paper
  defines
  the quantum indel distance
  and
  characterizes the insertion and deletion errors correcting capability of quantum codes
  by this distance. 
\end{itemize}

The rest of the paper is organized as follows.
Section \ref{sec:pre} introduces the notations used throughout the paper.
To efficiently handle insertion and deletion errors,
Section \ref{sec:ins_state} explicitly provides
quantum states after insertion error
and
Section \ref{sec:ins_del_state}
proves that
any composite errors consisting of insertion and deletion errors
are in
an error where $t$-insertion errors occur after $s$-deletion errors.
Section \ref{sub:relat_indel_correct} proves that
any quantum $t$-deletion-correcting code
is also
a quantum code correcting
a total of $t$ insertion and deletion errors,
whereas
any quantum $t$-insertion-correcting code
is not
a quantum $t$-deletion-correcting code.
Section \ref{sec:Qindel_dis} proposes the quantum indel distance and
describes insertion and deletion errors correcting capability of quantum codes
by this distance. 

%%%%%%%%%%%%%%%%%%%%%%%%%%%%%%%%%%%%%%%%%%%%%%%%%%%%%%%%%%%%%%%%%%%%%%%%%%%%%%%%%%%%%%%%%
\section{Preliminary}
\label{sec:pre}
This section gives the notations used throughout the paper.
Moreover,
this section introduces
quantum states, deletion errors, and insertion errors.

\subsection{Notations}
Let $\Z$, $\Z^+$, $\N$, $\NNR$, and $\C$ be
the sets of all integers, positive integers, nonnegative integers,
nonnegative real numbers, and complex numbers, respectively.
For $l \in \Z^+$,
define $\Z_l := \{ 0,1,...,l-1 \}$.
For a given set $S$,
let $|S|$ be the cardinality of $S$
and
$S^n$ be the set of all sequences of length $n$ over $S$.
For $n, m \in \mathbb{Z}$,
define
\begin{align*}
  \setint{m,n}
  &:=
  \{ i \in \Z \mid m \leq i \leq n \},\\
  \setint{n}
  &:=
  \setint{1,n},\\
  \subint{n}{m}
  &:=
  \{ S \subset \setint{n} \mid |S|=m \}.
\end{align*}
The complex conjugate and modulus of $c \in \C$ are written as $c^*$ and $|c|$,
respectively.

The Kronecker product of
an $m \times n$ matrix $\mat{M}_1=(m_{i,j})$ and a $p \times q$ matrix $\mat{M}_2$
is defined by
the $mp \times nq$ matrix
\begin{align*}
  \mat{M}_1 \otimes \mat{M}_2
  :=
  \begin{pmatrix}
    m_{1,1}\mat{M}_2 & m_{1,2}\mat{M}_2 & \cdots & m_{1,n}\mat{M}_2\\
    m_{2,1}\mat{M}_2 & m_{2,2}\mat{M}_2 & \cdots & m_{2,n}\mat{M}_2\\
    \vdots & \vdots & \ddots & \vdots\\
    m_{m,1}\mat{M}_2 & m_{m,2}\mat{M}_2 & \cdots & m_{m,n}\mat{M}_2\\
  \end{pmatrix}.
\end{align*}
We denote the set of all square matrices of order $n$ over $\C$, by $M(n)$.
For $\mat{M} \in M(n)$,
let $\text{tr}(\mat{M})$ be
the trace of $\mat{M}$
and
let $\mat{M}^{\dagger}$ be
the adjoint of $\mat{M}$.
Let $\mat{I}_n$ be
the identity matrix of order $n$.
A matrix $\mat{M} \in M(n)$ is unitary
if $\mat{M}^\dagger \mat{M} = \mat{I}_n$,
Hermitian
if $\mat{M}^\dagger = \mat{M}$.
A Hermitian matrix $\mat{M} \in M(n)$ is positive semidefinite
if $\bm{x}^\dagger \mat{M} \bm{x} \ge 0$
for all column vectors $\bm{x} \in \C^n$.
The notation $\mat{M} \ge 0$
means that
the Hermitian matrix $\mat{M}$
is
positive semidefinite.

Let $\pure{l}{n}$ be
the $l^n$-dimensional complex vector space.
For $x \in \mathbb{Z}_l$,
let $\ket{x} \in \C^l$ be a column vector
whose only the $(x+1)$th component is $1$ and the other components are $0$.
For $\bm{x} = (x_i) \in \mathbb{Z}_l^n$,
define
\begin{align*}
  \ket{\bm{x}}
  :=
  \ket{x_1} \otimes \ket{x_2} \otimes \cdots \otimes \ket{x_n} 
  \in
  \pure{l}{n}.
\end{align*}
Define
$\bra{\phi} := \ket{\phi}^{\dagger}$.
In this paper,
we denote an orthonormal basis for $\pure{l}{n}$
by $\{ \ket{ \bm{x}_L } \}_{ \bm{x} \in \Z_l^n }$.
Note that
$\ket{ \bm{x}_L }$ is not necessarily equal to $\ket{\bm{x}}$.

Given
any orthonormal basis $\{ \ket{ \bm{x}_L } \}_{ \bm{x} \in \Z_l^n }$ for $\pure{l}{n}$,
$\mat{M} \in M(l^n)$ is represented by
\begin{align}
  \mat{M}
  =
  \sum_{\bm{x},\bm{y} \in \Z_l^n}
  c_{\bm{x},\bm{y}}
  \ket{\bm{x}_L}\bra{\bm{y}_L}
  ,~~(c_{\bm{x}, \bm{y}} \in \C).
  \label{eq:orth_basis_xL}
\end{align}
In particular, 
\begin{align}
  \mat{M}
  =
  \sum_{\bm{x},\bm{y} \in \Z_l^n}
  c_{\bm{x},\bm{y}}
  \ket{\bm{x}}\bra{\bm{y}}
  ,~~(c_{\bm{x}, \bm{y}} \in \C).
  \label{eq:orth_basis_x}
\end{align}

\subsection{Quantum States \cite{lidar2013quantum}}
Let $\state{l}{n}$ be
the set of all density matrices of quantum states
represented by $n$ $l$-level qudits,
i.e.,
\begin{align*}
  \state{l}{n}
  :=
  \{
  \rho \in M(l^n)
  \mid
  \rho \ge 0,
  \text{tr}(\rho) = 1
  \}.
\end{align*}
For simplicity,
let $\state{l}{0} := \Big\{ \begin{pmatrix} 1 \end{pmatrix} \Big\}$.
Define $\state{l}{\N} := \bigcup_{n \in \N} \state{l}{n}$.

By the spectral decomposition,
$\rho \in \state{l}{n}$ is represented by
\begin{align}
  \rho
  &=
  \sum_{ \bm{x} \in \Z_l^n }
  p_{ \bm{x} } \ket{ \bm{x}_L } \bra{ \bm{x}_L },
  \label{eq:spec}
\end{align}
where
$\{ \ket{ \bm{x}_L } \}_{ \bm{x} \in \Z_l^n }$
is an orthonormal basis for $\pure{l}{n}$
and
$p_{ \bm{x} } \in \NNR$ satisfies
$\sum_{ \bm{x} \in \Z_l^n } p_{ \bm{x} } = 1$.
A quantum state $\rho$ is pure,
if $\rho$ has rank $1$.
Otherwise $\rho$ is mixed.
When $\rho \in \state{l}{n}$ is pure,
there exists a vector $\ket{\phi} \in \pure{l}{n}$
such that $\rho = \ket{\phi}\bra{\phi}$.
Then,
we denote $\rho$, by $\ket{\phi}$.

\subsection{Quantum Errors}\label{subsec:error}
We denote
the set of quantum states after a quantum error $\QU{E}$ to a quantum state $\rho$,
by $\QU{E}(\rho)$.

\subsubsection{Deletion Errors}
Deletion errors
loses
a qudit of a multi-qudit system.
For $\rho \in \state{l}{n}$ given in Eq.~\eqref{eq:orth_basis_x} and
$p \in \setint{n}$,
the partial trace $\text{Tr}_p$ is defined as 
\begin{align*}
  \text{Tr}_p(\rho)
  :=
  \sum_{ \bm{x}=(x_i), \bm{y}=(y_i) \in \Z_l^n } 
  c_{\bm{x},\bm{y}} 
  \text{tr}( \ket{x_p}\bra{y_p} )
  \ket{ \bm{x}_{ \neg p } }\bra{ \bm{y}_{ \neg p } },
\end{align*}
where
\begin{align*}
  \bm{x}_{ \neg p } := ( x_1, ..., x_{p-1}, x_{p+1}, ..., x_n ).
\end{align*}
Then,
for $\rho \in \state{l}{n}$ and
a set of deletion indices
$P = \{ p_1,p_2,...,p_s \} \in \subint{n}{s}$ with $p_1 < p_2 < \cdots < p_s$,
define $s$-deletion errors $\QU{D}_P$ as 
\begin{align*}
  \QU{D}_P(\rho)
  :=
  \text{Tr}_{p_1} \circ \text{Tr}_{p_2} \circ \cdots \circ \text{Tr}_{p_s}(\rho),
\end{align*}
where $f \circ g$ is the composition of the maps $f$ and $g$.
For $\rho \in \state{l}{n}$ and $s \leq n$,
define $s$-deletion errors $\QU{D}^s$ as
\begin{align*}
  \QU{D}^s(\rho)
  &:=
  \big\{
  \QU{D}_P(\rho)
  \mid
  P \in \subint{n}{s}
  \big\}.
\end{align*}

\subsubsection{Insertion Errors}
Insertion error adds a new qudit to a multi-qudit system.
For $\rho \in \state{l}{n}$ and $Q \in \subint{n+t}{t}$,
define
\begin{align*}
  \QU{I}_Q( \rho ) 
  := 
  \{ 
  \sigma \in \state{l}{(n+t)} 
  \mid 
  \QU{D}_Q(\sigma) = \rho 
  \}.
\end{align*}
Then,
$t$-insertion errors at $Q$
changes from $\rho$ to $\sigma \in \QU{I}_Q(\rho)$.
For $\rho \in \state{l}{n}$,
define $t$-insertion errors $\QU{I}^t$ as
\begin{align*}
  \QU{I}^t(\rho)
  &:=
  \bigcup_{Q \in \subint{n+t}{t}}
  \QU{I}_Q(\rho).
\end{align*}
For a given pure state $\rho = \ket{\phi}\bra{\phi}$,
$\sigma \in \QU{I}_Q(\ket{\phi}\bra{\phi})$
is expressed as follows \cite{sibayama_single_ins_def}:

\begin{theo}[\cite{sibayama_single_ins_def}]
  \label{theo:after_ins}
  Consider
  a quantum state $\rho \in \state{l}{n}$ given in Eq.~\eqref{eq:orth_basis_x}
  and
  a permutation $\tau$ on $\setint{n}$.
  With some abuse of notation,
  we define the {\it index permutation} $\tau(\rho) \in \state{l}{n}$ for $\rho$ by
  \begin{align*}
    \tau(\rho)
    :=
    \hspace{-3mm}
    \sum_{\substack{\bm{x}=(x_i) \in \mathbb{Z}_l^n\\\bm{y}=(y_i) \in \mathbb{Z}_l^n}}
    \hspace{-3mm}
    c_{\bm{x},\bm{y}}
    \ket{x_{\tau(1)} \cdots x_{\tau(n)}}\bra{y_{\tau(1)} \cdots y_{\tau(n)}}.
  \end{align*}
  In addition,
  for $t,n \in \mathbb{Z}^+$ and
  $Q = \{ q_1, q_2, ..., q_t \} \subset \setint{n+t}$ with $q_1 < q_2 < \cdots < q_t$,
  let $\tau^Q$ be the permutation on $\setint{n+t}$
  such that $\tau^Q(i)=q_i$ for $i \in \setint{n+1,n+t}$ and
  $j < k \Rightarrow \tau^Q(j) < \tau^Q(k)$ for $j,k \in \setint{n}$.
  Then,
  for a given pure state $\ket{\phi}\bra{\phi}$,
  any quantum state $\sigma \in \QU{I}_{Q}( \ket{\phi}\bra{\phi} )$ is denoted by
  \begin{align}
    \sigma
    =
    \tau^Q (\ket{\phi}\bra{\phi} \otimes \pi)
    \label{eq:pure_ins}
  \end{align}
  for some $\pi \in \state{l}{t}$.
\end{theo}

\subsubsection{Composite Errors}
Let $\QU{E}_2 \circ \QU{E}_1$ be
the error where a quantum error $\QU{E}_2$ occurs after a quantum error $\QU{E}_1$.
For example,
for $\rho \in \state{l}{n}$,
\begin{align*}
  \QU{D}^s \circ \QU{I}^t( \rho )
  &=
  \bigcup_{P \in \subint{n+t}{s}}
  \bigcup_{Q \in \subint{n+t}{t}}
  \{
  \QU{D}_P (\sigma)
  \mid
  \sigma \in \QU{I}_Q (\rho)
  \},\\
  \QU{I}^t \circ \QU{D}^s( \rho )
  &=
  \bigcup_{Q \in \subint{n-s+t}{t}}
  \bigcup_{\sigma \in \QU{D}^s (\rho)}
  \QU{I}_Q (\sigma).
\end{align*}
Moreover,
for $\QU{X} \subset \state{l}{\N}$,
define
\begin{align*}
  \QU{E}(\QU{X})
  :=
  \bigcup_{ \rho \in \QU{X} }
  \QU{E}(\rho).
\end{align*}
In addition,
we refer
an error where $s$-deletion errors and $t$-insertion errors occur compositely
as $(s,t)$-error.
For example,
$(1,2)$-error includes
\begin{align*}
  &\QU{D}^1 \circ \QU{I}^2,
  &&\QU{I}^1 \circ \QU{D}^1 \circ \QU{I}^1,
  &&\QU{I}^2 \circ \QU{D}^1.
\end{align*}
In general,
these are each distinct
because
deletion and insertion errors are non-commutative \cite{hagiwara2024introduction}.

%%%%%%%%%%%%%%%%%%%%%%%%%%%%%%%%%%%%%%%%%%%%%%%%%%%%%%%%%%%%%%%%%%%%%%%%%%%%%%%%%%%%%%%%%
\section{Quantum States after Insertion Errors to Mixed States}
\label{sec:ins_state}
This section describes
quantum states after insertion errors to mixed states
by generalizing Theorem \ref{theo:after_ins}.

A $k \times k$ principal submatrix of $\mat{M} \in M(n)$
is
a submatrix that lies on the same set of $k$ rows and columns,
and
a $k \times k$ principal minor is the determinant of a $k \times k$
principal submatrix.
\begin{exam}
  The $k \times k$ principal minors of
  $\begin{pmatrix} 1 & 2 & 3\\ 4 & 5 & 6\\ 7 & 8 & 9\end{pmatrix}$ are
    \begin{align*}
      &\begin{vmatrix}
        1 
      \end{vmatrix},
      \begin{vmatrix}
        5
      \end{vmatrix},
      \begin{vmatrix}
        9
      \end{vmatrix},
      &&(k=1)\\
      &\begin{vmatrix}
        1 & 2\\
        4 & 5
      \end{vmatrix},
      \begin{vmatrix}
        1 & 3\\
        7 & 9
      \end{vmatrix},
      \begin{vmatrix}
        5 & 6\\
        8 & 9
      \end{vmatrix},
      &&(k=2)\\
      &\begin{vmatrix}
        1 & 2 & 3\\
        4 & 5 & 6\\
        7 & 8 & 9
      \end{vmatrix}.
      &&(k=3)
    \end{align*}
\end{exam}

The following are  well known results for positive semidefinite matrices
(e.g., see \cite{horn1985matrix}).

\begin{lem}
  \label{lem:hansei}
  For a Hermitian matrix $\mat{M}$,
  $\mat{M} \ge 0$
  if and only if
  all principal minors of $\mat{M}$ are nonnegative.
\end{lem}

\begin{lem}
  \label{lem:sub_seitei_kai}
  For any matrix $\mat{A}$,
  $\mat{M} \ge 0$ satisfies
  $\mat{A} \mat{M} \mat{A}^\dagger \ge 0$.
\end{lem}

\begin{lem}
  \label{lem:px=0->cxy=0}
  For a positive semidefinite matrix $\mat{M} = (m_{i,j}) \in M(n)$,
  if $m_{i,i} = 0$,
  \begin{align*}
    \forall j \in \setint{n}
    ~~~
    m_{i,j} = m_{j,i} = 0
  \end{align*}
  holds.
\end{lem}

The following theorem provides
the quantum states after $\QU{I}_{\setint{n+1,n+t}}$.
\begin{theo}
  \label{theo:ins_state}
  Suppose
  $\rho \in \state{l}{n}$ is given by Eq.~\eqref{eq:spec}.
  Then,
  $\sigma \in \QU{I}_{\setint{n+1,n+t}}(\rho)$ is represented by
  \begin{align}
    \sigma
    =&
    \sum_{\bm{x},\bm{y} \in \Z_{l}^{n}}
    \sqrt{ p_{\bm{x}} p_{\bm{y}} }
    \ket{\bm{x}_L}\bra{\bm{y}_L} \otimes \mat{A}_{\bm{x},\bm{y}},
  \end{align}
  where
  $\mat{A}_{\bm{x},\bm{x}} \in \state{l}{t}$ and
  $\mat{A}_{\bm{x},\bm{y}} \in M(l^t)$ satisfies
  $\mat{A}_{\bm{x},\bm{y}}^\dagger = \mat{A}_{\bm{y},\bm{x}}$ and
  $\text{tr}( \mat{A}_{\bm{x},\bm{y}} ) = 0$ ($\bm{x} \neq \bm{y}$).
\end{theo}
\begin{proof}
  Since $\{ \ket{\bm{x}_L} \}_{ \bm{x} \in \Z_{l}^{n} }$ is   
  an orthonormal basis for $\pure{l}{n}$,
  $\{ \ket{\bm{x}_L} \otimes \ket{ \bm{z} } \}_{\bm{x} \in \Z_{l}^{n}, \bm{z} \in \Z_l^t}$
  is
  that for $\pure{l}{(n+t)}$.
  From Eq.~\eqref{eq:orth_basis_xL},
  $\sigma \in M(l^{(n+t)})$ is rewritten as
  \begin{align*}
    \sigma
    &=
    \sum_{\bm{x},\bm{y} \in \Z_{l}^{n}}
    \sum_{\bm{z}_1,\bm{z}_2 \in \Z_l^t}
    c_{\bm{x},\bm{z}_1,\bm{y},\bm{z}_2}
    (\ket{\bm{x}_L} \otimes \ket{ \bm{z}_1 })
    (\bra{\bm{y}_L} \otimes \bra{ \bm{z}_2 })\\
    &=
    \sum_{\bm{x},\bm{y} \in \Z_{l}^{n}}
    \ket{\bm{x}_L}\bra{\bm{y}_L} \otimes
    \mat{B}_{\bm{x},\bm{y}},
  \end{align*}
  where
  $\mat{B}_{\bm{x},\bm{y}} := \sum_{\bm{z}_1,\bm{z}_2 \in \Z_l^t} c_{\bm{x},\bm{z}_1,\bm{y},\bm{z}_2} \ket{ \bm{z}_1 }\bra{ \bm{z}_2 } \in M(l^t)$.
  Now,
  we will determine the conditions on $\mat{B}_{\bm{x},\bm{y}}$ from
  \begin{align*}
    \sigma
    =
    \sum_{\bm{x},\bm{y} \in \Z_{l}^{n}}
    \ket{\bm{x}_L}\bra{\bm{y}_L} \otimes \mat{B}_{\bm{x},\bm{y}}
    \in
    \QU{I}_{\setint{n+1,n+t}}(\rho).
  \end{align*}
  Since $\QU{D}_{\setint{n+1,n+t}}(\sigma) = \rho$,
  we get
  \begin{align}
    \text{tr} ( \mat{B}_{\bm{x},\bm{y}} )
    &=
    \begin{cases}
      p_{ \bm{x} } &( \text{if~} \bm{x} = \bm{y} )\\
      0 &( \text{if~} \bm{x} \neq \bm{y} )
    \end{cases}
    \label{eq:del_after}
  \end{align}
  Note that
  $\mat{U}' := \sum_{ \bm{x} \in \Z_l^n } \ket{\bm{x}} \bra{\bm{x}_L} \otimes \mat{I}_{l^t}$
  is unitary.
  From Lemma \ref{lem:sub_seitei_kai},
  $\sigma \ge 0$
  if and only if
  \begin{align*}
    \mat{U}' \sigma \mat{U}'^\dagger 
    =
    \sum_{\bm{x},\bm{y} \in \Z_{l}^{n}}
    \ket{\bm{x}}\bra{\bm{y}} \otimes
    \mat{B}_{\bm{x},\bm{y}}
    \ge 0.
  \end{align*}
  From Lemma \ref{lem:hansei},
  $\mat{B}_{\bm{x},\bm{x}} \ge 0$ and
  all diagonal elements of $\mat{B}_{\bm{x},\bm{x}}$ are nonnegative.
  From Eq.~\eqref{eq:del_after},
  all diagonal elements of $\mat{B}_{\bm{x},\bm{x}}$ are zero
  if $p_{\bm{x}} = 0$.
  Combining this and Lemma \ref{lem:px=0->cxy=0},
  if $p_{ \bm{x} } = 0$,
  $\mat{B}_{\bm{x},\bm{y}}$ and $\mat{B}_{\bm{y},\bm{x}}$ are equal to
  the zero matrix
  for any $\bm{y} \in \Z_q^n$.
  Hence,
  \begin{align*}
    \sigma
    &=
    \sum_{\bm{x},\bm{y} \in \Z_{l}^{n}}
    \sqrt{ p_{\bm{x}} p_{\bm{y}} }
    \ket{\bm{x}_L}\bra{\bm{y}_L} \otimes \mat{A}_{\bm{x},\bm{y}},
  \end{align*}
  where
  $\mat{A}_{\bm{x},\bm{y}} := \mat{B}_{\bm{x},\bm{y}} / \sqrt{ p_{\bm{x}} p_{\bm{y}} } \in M(l^t)$
  if
  $p_{\bm{x}} \neq 0$ and $p_{\bm{y}} \neq 0$.
  Since $\mat{B}_{\bm{x},\bm{x}} \ge 0$ and $\sigma$ is Hermitian,
  we get
  $\mat{A}_{\bm{x},\bm{x}} \ge 0$ and
  $\mat{A}_{\bm{x},\bm{y}}^\dagger = \mat{A}_{\bm{y},\bm{x}}$. 
  In addition,
  Eq.~\eqref{eq:del_after} yields
  \begin{align*}    
    \text{tr} ( \mat{A}_{\bm{x},\bm{y}} )
    &=
    \frac{1}{ \sqrt{ p_{\bm{x}} p_{\bm{y}} } }
    \text{tr}
    \left( \mat{B}_{\bm{x},\bm{y}} \right)
    =
    \begin{cases}
      1 & (\text{if~} \bm{x} = \bm{y})\\
      0 & (\text{if~} \bm{x} \neq \bm{y})
    \end{cases}.
  \end{align*}
  From the above,
  $\mat{A}_{\bm{x},\bm{x}} \in \state{l}{t}$ holds.
\end{proof}

By applying the index permutation to Theorem \ref{theo:ins_state},
any insertion errors for any insertion positions can be
expressed as follows:
\begin{coro}
  \label{coro:ins_state}
  Define $Q, \tau^Q(\rho)$ as in Theorem \ref{theo:after_ins}
  and
  $\mat{A}_{\bm{x},\bm{y}} \in M(l^t)$ as in Theorem \ref{theo:ins_state}.
  Then,
  for a quantum state $\rho \in \state{l}{n}$ in Eq.~\eqref{eq:spec},
  $\sigma \in \QU{I}_Q(\rho)$ is represented by
  \begin{align}
    \sigma
    =
    \tau^Q
    &
    \bigg(
    \sum_{\bm{x},\bm{y} \in \Z_{l}^{n}}
    (
    \sqrt{ p_{\bm{x}} p_{\bm{y}} }
    \ket{\bm{x}_L}\bra{\bm{y}_L} \otimes \mat{A}_{\bm{x},\bm{y}}
    )
    \bigg).
    \label{eq:theo:insstate_con_Q}
  \end{align}
\end{coro}

\begin{rem}
  If $\rho$ is pure,
  we get
  Eq.\eqref{eq:pure_ins} in Theorem \ref{theo:after_ins}
  from Eq.\eqref{eq:theo:insstate_con_Q} in Corollary \ref{coro:ins_state}.
\end{rem}

Noting that
Eq.~\eqref{eq:theo:insstate_con_Q} does not guarantee $\sigma \ge 0$,
we get from Corollary \ref{coro:ins_state} as follows:
\begin{coro}
  \label{coro:insset}
  Define $Q, \tau^Q(\rho)$ as in Theorem \ref{theo:after_ins}
  and
  $\rho \in \state{l}{n}$ as in Eq.~\eqref{eq:spec}.
  Then,
  \begin{align*}
    \QU{I}_Q(\rho)
    =
    \Bigg\{
    &\tau^Q
    \bigg(
    \sum_{\bm{x},\bm{y} \in \Z_{l}^{n}}
    (
    \sqrt{ p_{\bm{x}} p_{\bm{y}} }
    \ket{\bm{x}_L}\bra{\bm{y}_L} \otimes \mat{A}_{\bm{x},\bm{y}}
    )
    \bigg)
    \R{ \ge 0 }\\
    &\Big|
    \mat{A}_{\bm{x},\bm{y}} \in M(l^t),
    \text{tr}( \mat{A}_{\bm{x},\bm{y}} ) = 0 ~ (\bm{x} \neq \bm{y}),\\
    &~\mat{A}_{\bm{x},\bm{y}}^\dagger = \mat{A}_{\bm{y},\bm{x}},
    \mat{A}_{\bm{x},\bm{x}} \in \state{l}{t}
    \Bigg\}.
  \end{align*}
\end{coro}

\begin{exam}
  \label{exam:ins_set}
  For $p_0, p_1 \in \NNR$,
  define
  \begin{align*}
    \rho
    :=
    p_0\ket{00}\bra{00}
    +
    p_1\ket{11}\bra{11}
    \in
    \state{2}{2}.
  \end{align*}
  Then,
  from Corollary \ref{coro:ins_state},
  $\sigma_i \in \QU{I}_{\{i\}}(\rho)$ with $i=1,2,3$ are
  \begin{align*}
    \sigma_1
    =&
    p_0 ~ \pi_{00} \otimes \ket{00}\bra{00}
    +
    p_1 ~ \pi_{11} \otimes \ket{11}\bra{11}\\
    &+
    \sqrt{p_1 p_0} \mat{A} \otimes \ket{11}\bra{00}
    +
    \sqrt{p_0 p_1} \mat{A}^\dagger \otimes \ket{00}\bra{11},\\
    \sigma_2
    =&
    p_0 \ket{0}\bra{0} \otimes \pi_{00} \otimes \ket{0}\bra{0}
    +
    p_1 \ket{1}\bra{1} \otimes \pi_{11} \otimes \ket{1}\bra{1}\\
    &+
    \sqrt{p_1 p_0} \ket{1}\bra{0} \otimes \mat{A} \otimes \ket{1}\bra{0}\\
    &+
    \sqrt{p_0 p_1} \ket{0}\bra{1} \otimes \mat{A}^\dagger \otimes \ket{0}\bra{1},\\
    \sigma_3
    =&
    p_0 \ket{00}\bra{00} \otimes \pi_{00}
    +
    p_1 \ket{11}\bra{11} \otimes \pi_{11}\\
    &+
    \sqrt{p_1 p_0} \ket{11}\bra{00} \otimes \mat{A}
    +
    \sqrt{p_0 p_1} \ket{00}\bra{11} \otimes \mat{A}^\dagger,
  \end{align*}
  where
  $\pi_{00},\pi_{11} \in S(\C^2)$ and 
  $\mat{A} \in M(2)$ such that $\text{tr}(\mat{A}) = 0$.
\end{exam}

%%%%%%%%%%%%%%%%%%%%%%%%%%%%%%%%%%%%%%%%%%%%%%%%%%%%%%%%%%%%%%%%%%%%%%%%%%%%%%%%%%%
\section{Quantum States after Composite Error}
\label{sec:ins_del_state}
This section
proves that any $(s,t)$-error is included in $\QU{I}^t \circ \QU{D}^s$.
We give the following lemma for proving this.

\begin{lem}
  \label{lem:1ins_1del_kai}
  Any set $\QU{X} \subset \state{l}{n}$ satisfies
  \begin{align*}
    \QU{D}^1 \circ \QU{I}^1(\QU{X})
    \subset
    \QU{I}^1 \circ \QU{D}^1(\QU{X}).
  \end{align*}
\end{lem}

Before proving it,
we give an example.

\begin{exam}
  \label{exam:swap_TinsTdel}
  Define $\sigma_1$, $\sigma_2$, $\sigma_3$ and $\rho$ as in
  Example \ref{exam:ins_set}.
  Then,
  we get
  \begin{align*}
    \QU{D}_{\{1\}}(\sigma_1)
    =
    \QU{D}_{\{2\}}(\sigma_2)
    &=
    \QU{D}_{\{3\}}(\sigma_3)
    =
    \rho,\\
    \QU{D}_{\{2\}}(\sigma_1)
    =
    \QU{D}_{\{3\}}(\sigma_1)
    &=
    \QU{D}_{\{1\}}(\sigma_2)\\
    &=
    p_0 ~ \pi_{00} \otimes \ket{0}\bra{0}
    +
    p_1 ~ \pi_{11} \otimes \ket{1}\bra{1},\\
    \QU{D}_{\{3\}}(\sigma_2)
    =
    \QU{D}_{\{1\}}(\sigma_3)
    &=
    \QU{D}_{\{2\}}(\sigma_3)\\
    &=
    p_0 \ket{0}\bra{0} \otimes \pi_{00}
    +
    p_1 \ket{1}\bra{1} \otimes \pi_{11}.
  \end{align*}
  Hence,
  \begin{align*}
    &\QU{D}^1 \circ \QU{I}^1(\rho)\\
    &=
    \{ \rho \}\\
    &~~~\cup
    \left\{
    p_0 \pi_{00} \otimes \ket{0}\bra{0}
    +
    p_1 \pi_{11} \otimes \ket{1}\bra{1}
    \mid
    \pi_{00}, \pi_{11} \in S(\C^2)
    \right\}\\
    &~~~\cup
    \left\{
    p_0 \ket{0}\bra{0} \otimes \pi_{00}
    +
    p_1 \ket{1}\bra{1} \otimes \pi_{11}
    \mid
    \pi_{00}, \pi_{11} \in S(\C^2)
    \right\}
  \end{align*}
  Since
  $\QU{D}^1(\rho) = \{ p_0 \ket{0}\bra{0} + p_1 \ket{1}\bra{1} \}$,
  from Corollary \ref{coro:insset},
  we have
  \begin{align*}
    \QU{I}_{\{1\}} \circ &\QU{D}^1(\rho)\\
    =
    \Big\{
    &p_0 ~ \pi_{00} \otimes \ket{0}\bra{0}
    +p_1 ~ \pi_{11} \otimes \ket{1}\bra{1}\\
    &+ \sqrt{p_1p_0} \mat{A} \otimes \ket{1}\bra{0}
    + \sqrt{p_0p_1} \mat{A}^\dagger \otimes \ket{0}\bra{1} \ge 0\\
    &\mid
    \pi_{00}, \pi_{11} \in \state{2}{1},
    \mat{A} \in M(2),
    \text{tr}(\mat{A}) = 0
    \Big\},\\
    \QU{I}_{\{2\}} \circ &\QU{D}^1(\rho)\\
    =
    \Big\{
    &p_0 \ket{0}\bra{0} \otimes \pi_{00}
    +p_1 \ket{1}\bra{1} \otimes \pi_{11}\\
    &+ \sqrt{p_1p_0} \ket{1}\bra{0} \otimes \mat{A}
    + \sqrt{p_0p_1} \ket{0}\bra{1} \otimes \mat{A}^\dagger \ge 0\\
    &\mid
    \pi_{00}, \pi_{11} \in \state{2}{1},
    \mat{A} \in M(2),
    \text{tr}(\mat{A}) = 0
    \Big\}.
  \end{align*}
  From the above,
  $\QU{D}^1 \circ \QU{I}^1(\rho) \subset \QU{I}^1 \circ \QU{D}^1(\rho)$ holds.
  When $p_0, p_1 \neq 0$,
  $\QU{D}^1 \circ \QU{I}^1(\rho) \neq \QU{I}^1 \circ \QU{D}^1(\rho)$ holds
  since
  \begin{align*}
    \ket{\psi}
    :=
    \sqrt{p_0} 
    \ket{01}
    +
    \sqrt{p_1} 
    \ket{10}
    \in
    \QU{I}^1 \circ \QU{D}^1(\rho) \setminus \QU{D}^1 \circ \QU{I}^1(\rho).
  \end{align*} 
\end{exam}

\begin{rem}
  For some $\rho$,
  $\QU{D}^1 \circ \QU{I}^1(\rho) = \QU{I}^1 \circ \QU{D}^1(\rho)$ holds.
  For example,
  $\rho = \ket{00}\bra{00}$ satisfies
  \begin{align*}
    &\QU{D}^1 \circ \QU{I}^1(\rho)
    =
    \QU{I}^1 \circ \QU{D}^1(\rho)\\
    &=
    \bigcup_{\pi \in S(\C^2)}
    \{
    \ket{0}\bra{0} \otimes \pi,
    \pi \otimes \ket{0}\bra{0} \}.
  \end{align*}
\end{rem}

We give the following lemma to prove Lemma \ref{lem:1ins_1del_kai}.

\begin{lem}
  \label{lem_ID=non}
  For $P \in \subint{n}{t}$,
  $\QU{X} \subset \state{l}{n}$ satisfies
  \begin{align*}
    \QU{X}
    \subset
    \QU{I}_P \circ \QU{D}_P(\QU{X}).
  \end{align*} 
\end{lem}
\begin{proof}
  Since $\rho \in \QU{X} \subset \state{l}{n}$ satisfies
  \begin{align*}
    \QU{I}_P \circ \QU{D}_P(\{\rho\})
    &=
    \QU{I}_P( \QU{D}_P( \rho ) )\\
    &=
    \{
    \pi \in \state{l}{n}
    \mid
    \QU{D}_P(\pi) = \QU{D}_P(\rho)
    \}
    \ni
    \rho,
  \end{align*}
  we get
  $\QU{X} \subset \QU{I}_P \circ \QU{D}_P(\QU{X})$.
\end{proof}

\begin{proof}[Proof of Lemma \ref{lem:1ins_1del_kai}]
  To simplify the notation,
  we write $\QU{D}_{\{ p \}}$ as $\QU{D}_p$
  and
  $\QU{I}_{\{ q \}}$ as $\QU{I}_q$.
  It suffices to prove that
  for $\rho \in \state{l}{n}$,
  \begin{align}
    \forall
    p,q \in \setint{n+1}
    ~~~
    \QU{D}_q \circ \QU{I}_p( \rho )
    \subset
    \QU{I}^1 \circ \QU{D}^1( \rho ).
    \label{eq:juubunn}    
  \end{align}
  
  If $n = 1$,
  Eq.~\eqref{eq:juubunn} holds
  since
  $\QU{I}^1 \circ \QU{D}^1( \rho ) = \state{l}{1}$.
  Hence,
  consider the case for $n \ge 2$.
  If $p = q$,
  from Lemma \ref{lem_ID=non},
  Eq.~\eqref{eq:juubunn} follows
  \begin{align*}
    \QU{D}_p \circ \QU{I}_p( \rho )
    &=
    \{ \rho \}
    \subset
    \QU{I}_p \circ \QU{D}_p( \rho ). 
  \end{align*}
  For $p < q$,
  by substituting $\QU{X} = \QU{D}_q \circ \QU{I}_p( \rho )$ into Lemma \ref{lem_ID=non},
  we have 
  \begin{align*}
    \QU{D}_q \circ \QU{I}_p( \rho )
    &\subset
    \QU{I}_p \circ \QU{D}_p \circ \QU{D}_q \circ \QU{I}_p( \rho )\\
    &=
    \QU{I}_p \circ \QU{D}_{q-1} \circ \QU{D}_p \circ \QU{I}_p( \rho )\\
    &=
    \QU{I}_p \circ \QU{D}_{q-1}( \rho )\\
    &\subset
    \QU{I}^1 \circ \QU{D}^1( \rho ),
  \end{align*}
  where
  the first equality
  follows
  from $\QU{D}_p \circ \QU{D}_q = \QU{D}_{q-1} \circ \QU{D}_p$.
  Similarly,
  we get Eq.~\eqref{eq:juubunn}
  for $p > q$.
  Thus,
  we get Eq.~\eqref{eq:juubunn}.
\end{proof}

Lemma \ref{lem:1ins_1del_kai} is extended
to multiple deletion and insertion errors as follows.

\begin{lem}
  \label{lem:swap_TinsTdel_kai}
  For $s, t \in \N$,
  $\QU{X} \subset \state{l}{n}$ satisfies
  \begin{align*}
    \QU{D}^s \circ \QU{I}^t(\QU{X})
    \subset
    \QU{I}^t \circ \QU{D}^s(\QU{X}).
  \end{align*}
\end{lem}

\begin{proof}
  From Lemma \ref{lem:1ins_1del_kai},
  we get
  \begin{align*}
    \QU{D}^{s} \circ \QU{I}^{t}(\QU{X})
    &=
    \underbrace{
      \QU{D}^{1} \circ \cdots \circ \QU{D}^{1} \circ \B{\QU{D}^{1}}
    }_{s~\text{times}}
    \circ 
    \R{\QU{I}^{1}} \circ \QU{I}^{t-1} 
    (\QU{X})\\
    &\subset
    \QU{D}^{1} \circ \cdots \circ \QU{D}^{1} \circ \R{\QU{I}^{1}}
    \circ 
    \B{\QU{D}^{1}} \circ \QU{I}^{t-1} 
    (\QU{X})\\
    &\subset
    \QU{D}^{1} \circ \cdots \circ \R{\QU{I}^{1}} \circ \QU{D}^{1}
    \circ 
    \QU{D}^{1} \circ \QU{I}^{t-1} 
    (\QU{X})\\
    &~~~~~\vdots\\
    &\subset
    \R{\QU{I}^{1}} \circ \cdots \circ \QU{D}^{1} \circ \QU{D}^{1} \circ \QU{D}^{1}
    \circ 
    \QU{I}^{t-1} 
    (\QU{X})\\
    &=
    \R{\QU{I}^{1}} \circ \QU{D}^{s}
    \circ 
    \QU{I}^{t-1} 
    (\QU{X}).
  \end{align*}
  Similarly,
  \begin{align*}
    \QU{D}^{s} \circ \QU{I}^{t}(\QU{X})
    &\subset
    \QU{I}^{1} \circ \QU{D}^{s} \circ \QU{I}^{t-1} 
    (\QU{X})\\
    &\subset
    \QU{I}^{2} \circ \QU{D}^{s} \circ \QU{I}^{t-2}
    (\QU{X})\\
    &~~~~~\vdots\\
    &\subset
    \QU{I}^{t} \circ \QU{D}^{s}(\QU{X})
  \end{align*}
  holds.
\end{proof}

The following theorem
proves that any $(s,t)$-error is included in $\QU{I}^t \circ \QU{D}^s$.
In other words,
if a quantum code corrects $\QU{I}^t \circ \QU{D}^s$,
then it also corrects any $(s,t)$-errors.

\begin{theo}
  \label{theo:insdel_sphere}
  For $\QU{X} \subset \state{l}{n}$,
  let $\QU{A}$ be
  the set of quantum states after $(s,t)$-errors.
  Then,
  $\QU{A} \subset \QU{I}^t \circ \QU{D}^s(\QU{X})$ holds.
\end{theo}

\begin{proof}
  Any $(s,t)$-error is represented by
  \begin{align*}
    \QU{I}^{t_{1}} \circ \QU{D}^{s_{1}} \circ
    \QU{I}^{t_{2}} \circ \QU{D}^{s_{s}} \circ
    \cdots \circ
    \QU{I}^{t_{k}} \circ \QU{D}^{s_{k}}
  \end{align*}
  for some $k \in \Z^+$,
  $(s_i) \in \QU{S}_k$,
  $(t_i) \in \QU{T}_k$,
  where
  \begin{align*}
    \QU{S}_k
    &:=
    \left\{
    (s_i)
    \in
    (\Z^+)^{k-1}
    \times
    \N
    \mid
    s = \sum_{i=1}^k s_{i}
    \right\},\\
    \QU{T}_k
    &:=
    \left\{
    (t_i)
    \in
    \N
    \times
    (\Z^+)^{k-1}
    \mid
    t = \sum_{i=1}^k t_{i}
    \right\}.
  \end{align*} 
  Define $\QU{A} := \bigcup_{k \in \Z^+} \QU{A}_k$,
  where
  \begin{align*}
    \QU{A}_k
    &:=
    \hspace*{-2mm}
    \bigcup_{\bm{s} \in \QU{S}_k, \bm{t} \in \QU{T}_k}
    \hspace*{-2mm}
    \QU{I}^{t_{1}} \circ \QU{D}^{s_{1}} \circ
    \QU{I}^{t_{2}} \circ \QU{D}^{s_{2}} \circ
    \cdots \circ
    \QU{I}^{t_{k}} \circ \QU{D}^{s_{k}}
    (\QU{X}).
  \end{align*}    
  From Lemma \ref{lem:swap_TinsTdel_kai},
  we get
  \begin{align*}
    &\QU{A}_{k+1}\\
    &=
    \hspace*{-2mm}
    \bigcup_{\bm{s} \in \QU{S}_{k+1}, \bm{t} \in \QU{T}_{k+1}}
    \hspace*{-2mm}
    \QU{I}^{t_{1}} \circ \QU{D}^{s_{1}} \circ
    \cdots \circ
    \QU{I}^{t_{k}} \circ \B{\QU{D}^{s_{k}}} \circ
    \R{\QU{I}^{t_{k+1}}} \circ \QU{D}^{s_{k+1}}
    (\QU{X})\\
    &\subset
    \hspace*{-2mm}
    \bigcup_{\bm{s} \in \QU{S}_{k+1}, \bm{t} \in \QU{T}_{k+1}}
    \hspace*{-2mm}
    \QU{I}^{t_{1}} \circ \QU{D}^{s_{1}} \circ
    \cdots \circ
    \QU{I}^{t_{k}} \circ \R{\QU{I}^{t_{k+1}}} \circ
    \B{\QU{D}^{s_{k}}} \circ \QU{D}^{s_{k+1}}
    (\QU{X})\\
    &=
    \hspace*{-2mm}
    \bigcup_{\bm{s} \in \QU{S}_{k+1}, \bm{t} \in \QU{T}_{k+1}}
    \hspace*{-2mm}
    \QU{I}^{t_{1}} \circ \QU{D}^{s_{1}} \circ
    \cdots \circ
    \QU{I}^{t_{k}+t_{k+1}} \circ
    \QU{D}^{s_{k}+s_{k+1}}
    (\QU{X})\\
    &=
    \QU{A}_{k}
  \end{align*}
  Hence,
  we have
  \begin{align*}
    \cdots
    \subset \QU{A}_{k+1}
    \subset \QU{A}_{k}
    \subset \QU{A}_{k-1}
    \subset 
    \cdots    
    \subset \QU{A}_{1}
    =
    \QU{I}^t \circ \QU{D}^s(\QU{X}).
  \end{align*}
  Therefore,
  $\QU{A} = \bigcup_{k \in \Z^+} \QU{A}_k \subset \QU{I}^t \circ \QU{D}^s(\QU{X})$
  holds.
\end{proof}

%%%%%%%%%%%%%%%%%%%%%%%%%%%%%%%%%%%%%%%%%%%%%%%%%%%%%%%%%%%%%%%%%%%%%%%%%%%%%%
\section{Deletion Correcting Codes Correct Insertion Errors}
\label{sub:relat_indel_correct}
This paper employs the following definition of the error correcting capability
similar to classical coding theory.

\begin{defi}
  \label{def:error_corrct}
  For $\QU{X} \subset \state{l}{n}$ and a quantum error $\QU{E}$,
  we say that $\QU{X}$ corrects $\QU{E}$
  if $\QU{X}$ satisfies
  \begin{align*}
    \forall
    \rho_1, \rho_2
    \in
    \QU{X}
    ~
    \rho_1 \neq \rho_2
    \Rightarrow
    \QU{E}(\rho_1)
    \cap
    \QU{E}(\rho_2)
    =
    \emptyset.
  \end{align*}
\end{defi}

\subsection{Properties of Deletion and Insertion Errors}
\label{subsec:prop_indel}
This section derives the properties of deletion and insertion errors.
From the definition of $\QU{I}_Q(\rho_2)$,
we get
\begin{align}
  \rho_1 \in \QU{I}_Q(\rho_2)
  &\iff
  \QU{D}_Q(\rho_1) = \rho_2.
  \label{eq:Qins_def_iff}
\end{align}
The following lemma is
an extension of Eq.~\eqref{eq:Qins_def_iff}.

\begin{lem}
  \label{lem:sins_iff_sdel}
  Any $\rho_1 \in \state{l}{(m+s)}, \rho_2 \in \state{l}{m}$ satisfy
  \begin{align*}
    \rho_1 \in \QU{I}^s(\rho_2)
    \iff
    \QU{D}^s(\rho_1) \ni \rho_2.
  \end{align*}
\end{lem}

\begin{proof}
  From Eq.~\eqref{eq:Qins_def_iff},
  \begin{align*}
    \rho_1 \in \QU{I}^s(\rho_2)
    \iff&
    \exists Q \in \subint{m+s}{s}
    ~
    \rho_1 \in \QU{I}_Q(\rho_2)\\
    \iff&
    \exists Q \in \subint{m+s}{s}
    ~
    \QU{D}_Q(\rho_1) = \rho_2\\
    \iff&
    \QU{D}^s(\rho_1) \ni \rho_2
  \end{align*}
  holds.
\end{proof}

The following lemma is an extension of Lemma \ref{lem:sins_iff_sdel}.

\begin{lem}
  \label{lem:sins_iff_sdel_kai}
  Any $\QU{X}_1, \QU{X}_2 \subset \state{l}{\N}$ satisfy
  \begin{align*}
    \QU{X}_1 \cap \QU{I}^s(\QU{X}_2)
    \neq
    \emptyset
    \iff
    \QU{D}^s(\QU{X}_1) \cap \QU{X}_2
    \neq
    \emptyset.
  \end{align*}
\end{lem}

\begin{proof}
  From Lemma \ref{lem:sins_iff_sdel},
  \begin{align*}
    \QU{X}_1 \cap \QU{I}^s(\QU{X}_2)
    \neq
    \emptyset
    \iff&
    \exists
    \rho_1 \in \QU{X}_1
    ~
    \rho_1 \in \QU{I}^s(\QU{X}_2)\\
    \iff&
    \exists
    \rho_1 \in \QU{X}_1
    ~
    \exists
    \rho_2 \in \QU{X}_2
    ~
    \rho_1 \in \QU{I}^s(\rho_2)\\
    \iff&
    \exists
    \rho_1 \in \QU{X}_1
    ~
    \exists
    \rho_2 \in \QU{X}_2
    ~
    \QU{D}^s(\rho_1) \ni \rho_2\\
    \iff&
    \exists
    \rho_2 \in \QU{X}_2
    ~
    \QU{D}^s(\QU{X}_1) \ni \rho_2\\
    \iff&
    \QU{D}^s(\QU{X}_1) \cap \QU{X}_2
    \neq
    \emptyset
  \end{align*}
  holds.
\end{proof}

The following is derived
from Lemma \ref{lem:sins_iff_sdel_kai}.

\begin{lem}
  \label{lem:sub_set}
  Any $\QU{X}_1,\QU{X}_2 \subset \state{l}{\N}$ satisfy
  \begin{align*}
    \QU{X}_1
    \cap
    \QU{I}^s \circ \QU{D}^t(\QU{X}_2)
    \neq
    \emptyset
    &\iff
    \QU{D}^s(\QU{X}_1) \cap \QU{D}^t(\QU{X}_2)
    \neq
    \emptyset,\\
    \QU{X}_1
    \cap
    \QU{D}^s \circ \QU{I}^t(\QU{X}_2)
    \neq
    \emptyset
    &\iff
    \QU{I}^s(\QU{X}_1) \cap \QU{I}^t(\QU{X}_2)
    \neq
    \emptyset.
  \end{align*}
  In particular,
  if $\QU{X}_1 = \{ \rho \}$,
  \begin{align*}
    \rho
    \in
    \QU{I}^s \circ \QU{D}^t(\QU{X}_2)
    &\iff
    \QU{D}^s(\rho) \cap \QU{D}^t(\QU{X}_2)
    \neq
    \emptyset,\\
    \rho
    \in
    \QU{D}^s \circ \QU{I}^t(\QU{X}_2)
    &\iff
    \QU{I}^s(\rho) \cap \QU{I}^t(\QU{X}_2)
    \neq
    \emptyset
  \end{align*}
  holds.
\end{lem}

\subsection{Main Result}
\begin{theo}
  \label{theo:1del_iff_1ins_2}
  If $\QU{X} \subset \state{l}{n}$ corrects $t$-deletion errors,
  then $\QU{X}$ also corrects a total of $t$ deletion and insertion errors.
\end{theo}

Theorem \ref{theo:1del_iff_1ins_2}
is derived
from the following Lemma.

\begin{lem}
  \label{lem:1del_iff_1ins_2}
  For quantum states $\rho_1, \rho_2 \in \state{l}{n}$,
  if $\QU{I}^{t} \circ \QU{D}^{s}(\rho_1) \cap \QU{I}^{t} \circ \QU{D}^{s}(\rho_2) =\emptyset$,
  then
  \begin{align*}
    \QU{I}^{(t+1)} \circ \QU{D}^{(s-1)}(\rho_1)
    \cap
    \QU{I}^{(t+1)} \circ \QU{D}^{(s-1)}(\rho_2)
    =
    \emptyset.
  \end{align*}
\end{lem}

\begin{proof}
  From Lemmas \ref{lem:sins_iff_sdel_kai} and \ref{lem:sub_set},
  \begin{align}
    &\QU{I}^{t} \circ \QU{D}^{s}(\rho_1)  \cap \QU{I}^{t} \circ \QU{D}^{s}(\rho_2)
    =
    \emptyset
    \nonumber\\
    \iff&
    \QU{D}^{s}(\rho_1)  \cap \QU{D}^{t} \circ \QU{I}^{t} \circ \QU{D}^{s}(\rho_2)
    =
    \emptyset
    \nonumber\\
    \iff&
    \rho_1
    \notin
    \QU{I}^{s} \circ \QU{D}^{t} \circ \QU{I}^{t} \circ \QU{D}^{s}(\rho_2)
    \label{eq:assu_change1}
  \end{align}
  holds.
  From Lemma \ref{lem:swap_TinsTdel_kai},
  we get
  \begin{align*}
    &\QU{I}^{(s-1)}\circ \QU{D}^{(t+1)} \circ \QU{I}^{(t+1)} \circ \QU{D}^{(s-1)}(\rho_2)\\
    &\subset
    \QU{I}^{s} \circ \QU{D}^{(t+1)} \circ \QU{I}^{t} \circ \QU{D}^{(s-1)}(\rho_2)\\
    &\subset 
    \QU{I}^{s} \circ \QU{D}^{t} \circ \QU{I}^{t} \circ \QU{D}^{s}(\rho_2).
  \end{align*}
  Hence,
  we have
  \begin{align}
    &\rho_1 \notin \QU{I}^{s} \circ \QU{D}^{t} \circ \QU{I}^{t} \circ \QU{D}^{s}(\rho_2)
    \nonumber\\
    &\Rightarrow
    \rho_1 \notin \QU{I}^{(s-1)}\circ \QU{D}^{(t+1)} \circ \QU{I}^{(t+1)} \circ \QU{D}^{(s-1)}(\rho_2).
    \label{eq:st->s-1t+1}
  \end{align}
  By a similar way to Eq.~\eqref{eq:assu_change1},
  \begin{align}
    &\rho_1
    \notin
    \QU{I}^{(s-1)} \circ \QU{D}^{(t+1)} \circ \QU{I}^{(t+1)} \circ \QU{D}^{(s-1)}(\rho_2)
    \nonumber\\
    \iff
    &\QU{I}^{(t+1)} \circ \QU{D}^{(s-1)}(\rho_1)
    \cap
    \QU{I}^{(t+1)} \circ \QU{D}^{(s-1)}(\rho_2)
    =
    \emptyset
    \label{eq:assu_change2}
  \end{align}
  holds.
  Combining
  Eqs.~\eqref{eq:assu_change1}, \eqref{eq:st->s-1t+1}, and \eqref{eq:assu_change2},
  we get the lemma.
\end{proof}

Since deletion and insertion errors are non-commutative \cite{hagiwara2024introduction},
the converse of Theorem \ref{theo:1del_iff_1ins_2} does not necessarily hold.

\begin{exam}
  Define
  $\rho$ and $\ket{\psi}$ as in Example \ref{exam:ins_set} and \ref{exam:swap_TinsTdel},
  respectively.
  Then,
  we get
  \begin{align*}
    &\ket{\psi}\bra{\psi}
    \in
    \QU{I}^1 \circ \QU{D}^1(\rho),
    &&\ket{\psi}\bra{\psi}
    \notin
    \QU{D}^1 \circ \QU{I}^1(\rho).
  \end{align*}
  From Lemma \ref{lem:sub_set},
  we get
  \begin{align*}
    &\QU{D}^1(\ket{\psi}\bra{\psi}) \cap \QU{D}^1(\rho)
    \neq
    \emptyset,
    &&\QU{I}^1(\ket{\psi}\bra{\psi}) \cap \QU{I}^1(\rho)
    =
    \emptyset.
  \end{align*}
  In other words,
  $\{ \rho,\ket{\psi}\bra{\psi} \}$ corrects a single insertion
  but not a single deletion.
\end{exam}

%%%%%%%%%%%%%%%%%%%%%%%%%%%%%%%%%%%%%%%%%%%%%%%%%%%%%%%%%%%%%%%%%%%%%%%%%%%%%%
\section{Quantum Indel Distance}
\label{sec:Qindel_dis}
This section defines the quantum indel distance and
describes 
insertion and deletion errors correcting capability of quantum codes
by this distance.

\subsection{Definition and Examples}
The indel distance for $\bm{x} \in \Z_l^n$ and $\bm{y} \in \Z_l^m$ is
defined as the sum of the number of classical deletions and insertions
required to transform $\bm{x}$ into $\bm{y}$ \cite{levenshtein1966binary}.
Similar to the indel distance,
we define
the quantum indel distance for $\rho_1, \rho_2 \in \state{l}{\N}$ as
the sum of the number of deletion and insertion errors
required to transform
$\rho_1$ into $\rho_2$.
The quantum indel distance is mathematically defined as follows:

\begin{defi}
  The quantum indel distance $d( \rho_1, \rho_2 )$ of $\rho_1, \rho_2 \in \state{l}{\N}$
  is defined as follows.
  \begin{align*}
    d( \rho_1, \rho_2 )
    &:=
    \min_{(s,t) \in \N^2, (s,t)\text{-error~}\QU{E}}
    \{
    s+t
    \mid
    \rho_2
    \in
    \QU{E}(\rho_1)
    \}\\
    &=
    \min_{(s,t) \in \N^2}
    \{
    s+t
    \mid
    \rho_2
    \in
    \QU{I}^t \circ \QU{D}^s(\rho_1)
    \}\\
    &=
    \min_{(s,t) \in \N^2}
    \{
    s+t
    \mid
    \QU{D}^s(\rho_1) \cap \QU{D}^t(\rho_2)
    \neq
    \emptyset
    \}
  \end{align*}
\end{defi}

The first equality
follows
from Theorem \ref{theo:insdel_sphere}
and
the second equality
follows
from Lemma \ref{lem:sub_set}.

\begin{exam}
  Suppose
  \begin{align*}
    \rho_1
    &=
    \tfrac{1}{2}
    \ket{00}\bra{00}
    +
    \tfrac{1}{2}
    \ket{11}\bra{11},\\
    \rho_2
    &=
    \tfrac{1}{2}
    \ket{10}\bra{10}
    +
    \tfrac{1}{2}
    \ket{01}\bra{01}.
  \end{align*}
  Then,
  $d( \rho_1, \rho_2 ) = 2$ holds
  since
  \begin{align*}
    \QU{D}_{\{ 1 \}}(\rho_1)
    =
    \QU{D}_{\{ 2 \}}(\rho_2)
    =
    \tfrac{1}{2}\ket{0}\bra{0}
    +
    \tfrac{1}{2}\ket{1}\bra{1}.
  \end{align*} 
\end{exam}

\begin{defi}
  The minimum distance $d_{\text{min}}(\QU{X})$ of $\QU{X} \subset \state{l}{\N}$
  is defined as follows:
  \begin{align*}
    d_{\text{min}}(\QU{X})
    :=
    \min_{ \rho_1, \rho_2 \in \QU{X} : \rho_1 \neq \rho_2 }
    d( \rho_1, \rho_2 )
  \end{align*}
\end{defi}

Now,
we will calculate the minimum distances of several quantum codes.
Let $j$ be the imaginary unit.

\begin{exam}
  Define
  $\QU{X}_1 := \{ \alpha \ket{00} + \beta \ket{11} \mid \alpha,\beta \in \C, |\alpha|^2+|\beta|^2 = 1 \}$.
  Note that
  \begin{align*}
    \rho
    &=
    |\alpha|^2 \ket{00}\bra{00} + \alpha\beta^* \ket{00}\bra{11}\\
    &~~~+
    \beta\alpha^* \ket{11}\bra{00} + |\beta|^2 \ket{11}\bra{11}
    \in \QU{X}_1
  \end{align*}
  We have
  \begin{align*}
    \QU{D}_{\{1\}}(\rho)
    =
    \QU{D}_{\{2\}}(\rho)
    =
    |\alpha|^2 \ket{0}\bra{0} + |\beta|^2 \ket{1}\bra{1}.
  \end{align*}
  Hence,
  for $\theta \in \NNR$,
  \begin{align*}
    \ket{\psi_1}
    &:=
    \alpha \ket{00} + \beta \ket{11} \in \QU{X}_1,\\
    \ket{\psi_2}
    &:=
    \alpha \ket{00} + e^{j \theta}\beta \ket{11} \in \QU{X}_1
  \end{align*}
  satisfy
  \begin{align*}
    |\alpha|^2 \ket{0}\bra{0} + |\beta|^2 \ket{1}\bra{1}
    \in
    \QU{D}^1(\ket{\psi_1}\bra{\psi_1})
    \cap
    \QU{D}^1(\ket{\psi_2}\bra{\psi_2}).
  \end{align*}
  Thus, $d_{\text{min}}(\QU{X}_1) = 2$.
\end{exam}

\begin{exam}
  \label{exam:dis=4}
  Denote the set of $\bm{x} \in \{ 0,1 \}^n$ with Hamming weight $i$,
  by $X_{n,i}$.
  Define $\ket{i_{(n)}} := \sum_{ \bm{x} \in X_{n,i} } \ket{\bm{x}}$.
  Consider
  the minimum distance of the following quantum single-deletion-correcting code
  $\QU{X}_2$ given in \cite{hagiwara2020four}:
  \begin{align*}
    &\QU{X}_2
    :=
    \left\{
    \alpha
    \ket{0_L}
    +
    \beta
    \ket{1_L}
    \mid
    \alpha,\beta \in \C, |\alpha|^2+|\beta|^2 = 1
    \right\},
  \end{align*}
  where
  \begin{align*}
    &\ket{0_L}
    :=
    \tfrac{ 1 }{ \sqrt{2} }\big( \ket{0_{(4)}} + \ket{4_{(4)}}\big),
    ~~~~
    \ket{1_L}
    :=
    \tfrac{ 1 }{ \sqrt{6} } \ket{2_{(4)}}.
  \end{align*}
  For $p \in \setint{4}, \rho = \alpha \ket{0_L} + \beta \ket{1_L} \in \QU{X}_2$,
  we have
  \begin{align}
    &\QU{D}_{\{p\}}(\rho)
    \nonumber\\
    &=
    \tfrac{1}{2}
    \big( \alpha\ket{0_{(3)}} + \tfrac{\beta}{\sqrt{3}}\ket{2_{(3)}} \big)
    \big( \alpha^*\bra{0_{(3)}} + \tfrac{\beta^*}{\sqrt{3}}\bra{2_{(3)}} \big)
    \nonumber\\
    &~~~+
    \tfrac{1}{2}
    \big( \alpha\ket{3_{(3)}}+\tfrac{\beta}{\sqrt{3}}\ket{1_{(3)}} \big)
    \big( \alpha^*\bra{3_{(3)}}+\tfrac{\beta^*}{\sqrt{3}}\bra{1_{(3)}} \big).
    \label{eq:expm_del_code_dis}
  \end{align}
  For distinct $\rho_1, \rho_2 \in \QU{X}_2$,
  $\QU{D}^1(\rho_1) \cap \QU{D}^1(\rho_2) = \emptyset$
  since
  $\ket{0_{(3)}}, \ket{1_{(3)}}, \ket{2_{(3)}}, \ket{3_{(3)}}$
  are
  mutually orthogonal.
  That is,
  $d_{\text{min}}(\QU{X}_2) \ge 3$.
  From Eq.~\eqref{eq:expm_del_code_dis},
  for $P \in \subint{4}{2}$,
  we have
  \begin{align}
    \QU{D}_P(\rho)
    =&
    \tfrac{1}{2}
    \big(
    |\alpha|^2 + \tfrac{|\beta|^2}{3}
    \big)
    \big(
    \ket{0_{(2)}}\bra{0_{(2)}} + \ket{2_{(2)}}\bra{2_{(2)}}
    \big)
    \nonumber\\
    &+
    \tfrac{ \alpha\beta^* + \alpha^*\beta }{2\sqrt{3}}
    \big(
    \ket{0_{(2)}}\bra{2_{(2)}} + \ket{2_{(2)}}\bra{0_{(2)}}
    \big)
    \nonumber\\
    &+
    \tfrac{|\beta|^2}{3}\ket{1_{(2)}}\bra{1_{(2)}}.
    \label{eq:expm_del_code_dis_4}
  \end{align}
  Let $\text{arg } c$ be the argument of $c$.
  Then,
  \begin{align*}
    \ket{\psi_1}
    &:=
    \alpha \ket{0_L} + \beta \ket{1_L} \in \QU{X}_2,\\
    \ket{\psi_2}
    &:=
    \alpha \ket{0_L} + e^{2j( \text{arg } \alpha - \text{arg } \beta )}
    \beta \ket{1_L} \in \QU{X}_2
  \end{align*}
  satisfy
  \begin{align*}
    \text{(R.H.S. of Eq.~\eqref{eq:expm_del_code_dis_4})}
    \in
    \QU{D}^2(\ket{\psi_1}\bra{\psi_1})
    \cap
    \QU{D}^2(\ket{\psi_2}\bra{\psi_2}).
  \end{align*}
  Hence,
  $d_{\text{min}}(\QU{X}_2) = 4$.
\end{exam}

\subsection{Properties}
\begin{theo}
  The quantum indel distance is a metric on $\state{l}{\N}$.
\end{theo}

\begin{proof}
  Any $\rho_1, \rho_2 \in \state{l}{\N}$ satisfy
  \begin{align*}
    d(\rho_1, \rho_2) = 0
    &\iff
    \emptyset
    \neq
    \QU{D}^0(\rho_1) \cap \QU{D}^0(\rho_2)
    =
    \{ \rho_1 \} \cap \{ \rho_2 \}\\
    &\iff
    \rho_1 = \rho_2.
  \end{align*}
  The non-negativity and the symmetry are trivial.
  Now,
  we will prove the triangle inequality.
  For $\rho_1, \rho_2, \rho_3 \in \state{l}{\N}$,
  define
  \begin{align*}
    (s_1,t_1)
    &:=
    \underset{(s,t) \in \N^2}{\operatorname{argmin}}~
    \{
    s+t
    \mid
    \rho_2 \in \QU{I}^{t} \circ \QU{D}^{s}(\rho_1)
    \},\\
    (s_2,t_2)
    &:=
    \underset{(s,t) \in \N^2}{\operatorname{argmin}}~
    \{
    s+t
    \mid
    \rho_2 \in \QU{I}^{t} \circ \QU{D}^{s}(\rho_3)
    \}.
  \end{align*}
  From $\rho_2 \in \QU{I}^{t_1} \circ \QU{D}^{s_1}(\rho_1)$ and
  $\rho_2 \in \QU{I}^{t_2} \circ \QU{D}^{s_2}(\rho_3)$,
  we get
  \begin{align*}
    \QU{I}^{t_1} \circ \QU{D}^{s_1}(\rho_1)
    \cap
    \QU{I}^{t_2} \circ \QU{D}^{s_2}(\rho_3)
    \neq
    \emptyset.
  \end{align*}
  From Lemma \ref{lem:sub_set} and Theorem \ref{theo:insdel_sphere},
  we have
  \begin{align*}
    \rho_1 
    &\in
    \QU{I}^{s_1} \circ \QU{D}^{t_1} \circ \QU{I}^{t_2} \circ \QU{D}^{s_2}(\rho_3)
    \subset
    \QU{I}^{s_1+t_2} \circ \QU{D}^{t_1+s_2}(\rho_3).    
  \end{align*}
  Therefore,
  \begin{align*}
    d(\rho_1 ,\rho_3)
    &\leq
    s_1+t_2 + t_1+s_2\\
    &=
    d(\rho_1, \rho_2) + d(\rho_2, \rho_3).
  \end{align*}
  holds.
\end{proof}

\begin{lem}
  \label{lem:1del_iff_1ins_1}
  Any $\rho_1, \rho_2 \in \state{l}{n}$ satisfy
  \begin{align*}
    d(\rho_1,\rho_2)
    \ge
    2t+1
    \iff
    \QU{D}^t(\rho_1) \cap \QU{D}^t(\rho_2) = \emptyset
  \end{align*}
\end{lem}
Noting that
$d(\rho_1,\rho_2)$
is
even
for $\rho_1,\rho_2 \in \state{l}{n}$,
we get
Lemma \ref{lem:1del_iff_1ins_1}.
Lemma \ref{lem:1del_iff_1ins_1} yields the following theorem.

\begin{theo}
  \label{theo:1del_iff_1ins_1}
  The set $\QU{X} \subset \state{l}{n}$
  corrects
  $t$-deletion errors
  iff $d_{\text{min}}(\QU{X}) \ge 2t+1$.
\end{theo}

From Theorems \ref{theo:1del_iff_1ins_2} and \ref{theo:1del_iff_1ins_1},
we obtain the following.

\begin{coro}
  The set $\QU{X} \subset \state{l}{n}$
  corrects
  a total of $t$ insertion and deletion errors
  if $d_{\text{min}}(\QU{X}) \ge 2t+1$.
\end{coro}

\begin{exam}
  From Example \ref{exam:dis=4},
  $d_{\text{min}}(\QU{X}_2) =4$ holds.
  Hence,
  $\QU{X}_2$
  corrects
  a single insertion error,
  as shown in \cite{hagiwara2021four}.
\end{exam}

\section*{Acknowledgment}
This work was supported by USPS KAKENHI Grant Number 22K11905.

%\bibliographystyle{ieicetr}
\bibliography{myref}

%% \profile{Ken NAKAMURA}{
%%   received B.E. and M.E. degrees from Yamaguchi University
%%   in 2023 and 2025, respectively.
%%   Since April 2025,
%%   he has been currently pursuing the D.S.\ degree in Department of Informatics,
%%   Yamaguchi University.
%%   His research interests is insertion/deletion correcting quantum codes.
%%   He is a member of IEEE.
%% }

%% \profile{Takayuki Nozaki}{
%%   received B.E., M.E., and D.E.\ degrees from Tokyo Institute of Technology
%%   in 2008, 2010, and 2012, respectively.
%%   %% From 2010 to 2013, he was a Research Fellow of the JSPS.
%%   %% From 2013 to 2015, he was a Research Associate at Kanagawa University.
%%   %% From 2015 to 2017, he was an Assistant Professor at Yamaguchi University.
%%   %% Since 2017, he has been a Senior Lecturer at Yamaguchi University.
%%   He is currently an associate professor at Yamaguchi University.
%%   His research interests include error correcting codes and coding theory.
%%   He received the Best Paper Award from the IEICE of Japan in 2018.
%%   He is a member of IEEE, ACM, AMS, IPSJ, and JSTS.
%% }
\end{document}

